\section{USAJMO 2014}
        \begin{enumerate}
        	\item (\textit{USAJMO 2014}) Let $a$, $b$, $c$ be real numbers greater than or equal to $1$. Prove that 
 \[\min{\left (\frac{10a^2-5a+1}{b^2-5b+10},\frac{10b^2-5b+1}{c^2-5c+10},\frac{10c^2-5c+1}{a^2-5a+10}\right )}\leq abc.\]


 

	\item (\textit{USAJMO 2014}) Let $\triangle{ABC}$ be a non-equilateral, acute triangle with $\angle A=60^\circ$, and let $O$ and $H$ denote the circumcenter and orthocenter of $\triangle{ABC}$, respectively.


 (a) Prove that line $OH$ intersects both segments $AB$ and $AC$.


 (b) Line $OH$ intersects segments $AB$ and $AC$ at $P$ and $Q$, respectively. Denote by $s$ and $t$ the respective areas of triangle $APQ$ and quadrilateral $BPQC$. Determine the range of possible values for $s/t$.


 

	\item (\textit{USAJMO 2014}) Let $\mathbb{Z}$ be the set of integers. Find all functions $f : \mathbb{Z} \rightarrow \mathbb{Z}$ such that 
 \[xf(2f(y)-x)+y^2f(2x-f(y))=\frac{f(x)^2}{x}+f(yf(y))\] for all $x, y \in \mathbb{Z}$ with $x \neq 0$.


 

	\item (\textit{USAJMO 2014}) Let $b\geq 2$ be an integer, and let $s_b(n)$ denote the sum of the digits of $n$ when it is written in base $b$.  Show that there are infinitely many positive integers that cannot be represented in the form $n+s_b(n)$, where $n$ is a positive integer.


 

	\item (\textit{USAJMO 2014}) Let $k$ be a positive integer. Two players $A$ and $B$ play a game on an infinite grid of regular hexagons. Initially all the grid cells are empty. Then the players alternately take turns with $A$ moving first. In his move, $A$ may choose two adjacent hexagons in the grid which are empty and place a counter in both of them. In his move, $B$ may choose any counter on the board and remove it. If at any time there are $k$ consecutive grid cells in a line all of which contain a counter, $A$ wins. Find the minimum value of $k$ for which $A$ cannot win in a finite number of moves, or prove that no such minimum value exists.


 

	\item (\textit{USAJMO 2014}) Let $ABC$ be a triangle with incenter $I$, incircle $\gamma$ and circumcircle $\Gamma$.  Let $M,N,P$ be the midpoints of sides $\overline{BC}$, $\overline{CA}$, $\overline{AB}$ and let $E,F$ be the tangency points of $\gamma$ with $\overline{CA}$ and $\overline{AB}$, respectively.  Let $U,V$ be the intersections of line $EF$ with line $MN$ and line $MP$, respectively, and let $X$ be the midpoint of arc $BAC$ of $\Gamma$.


 (a) Prove that $I$ lies on ray $CV$.


 (b) Prove that line $XI$ bisects $\overline{UV}$.


 

\end{enumerate}